% inputs/layout.tex
% $Author: predrag $ $Date: 2021-07-25 19:52:46 -0400 (Sun, 25 Jul 2021) $

% from ChaosBook                        Predrag          5jun2008

\newcommand{\authorXD}[1]
     {\hfill (X. Ding, #1)}
\newcommand{\authorSFPC}[1] % Date 6 Jun 2009, for example
     {\hfill (S. Froehlich and P. Cvitanovi\'c, #1)}
\newcommand{\authorSF}[1]
     {\hfill (S. Froehlich, #1)}
\newcommand{\authorYL}[1]
     {\hfill (Y. Lan, #1)}
\newcommand{\authorHL}[1]
     {\hfill (H. Liang, #1)}
\newcommand{\authorRW}[1]
     {\hfill (R. Wilczek, #1)}
\newcommand{\authorES}[1]
     {\hfill (E. Siminos, #1)}
\newcommand{\authorPC}[1]
     {\hfill (P. Cvitanovi\'c, #1)}
\newcommand{\authorJMPC}
     {\hfill (J. Mathiesen and P. Cvitanovi\'c)}
\newcommand{\authorRA}%{\marginpar{RM+PC}}
     {\hfill (R. Artuso)}

        \ifdasbuch
%%%%%%%%%%%%%%% FORMATING: PAGINATION FOR bin/mkbook %%%%%%%%%%%%%%%%%%%%%%%
\newcounter{startSect}  %% page numbering used to generate ps pages
\newcounter{endSect}    %% page numbering used to generate ps pages
\newcommand{\chapName}{Unnamed} %% chapter label
\newcommand{\Chapter}[3]{\clearpage\thispagestyle{empty}\cleardoublepage
                % was \newcommand{\inputfile}{#1 -  #2} % until cludge 2021-05-04
                % PC 2018-01-28, get rid of it in ChaosBook?
                \renewcommand{\inputfile}{#1 -  #2}
                        \renewcommand{\chapName}{#1}
                         \chapter[#3]{#3}
                         \label{c-#1}
                         \setcounter{startSect}{\thepage}
                        }
\newcommand{\ChapterEnd}{\setcounter{endSect}{\thepage}
%               Predrag         13nov2010
                        \addtocontents{toc}{\par} % end ... references entry
                         }
%%%%%%%% for pointers:
\newcommand{\tangent}{\includegraphics[width=0.08\textwidth]{Fig/chili.eps}~~}
\newcommand{\toChap}[1]
{\marginpar[{\color{green}\sf\footnotesize\hfill chapter \ref{#1} }]%
{\color{green}\sf\footnotesize chapter \ref{#1} \hfill}}
\newcommand{\toAppe}[1]
{\marginpar[{\color{green}\sf\footnotesize\hfill appendix \ref{#1} }]
               {\color{green}\sf\footnotesize appendix \ref{#1} \hfill}}
\newcommand{\toSect}[1]
{\marginpar[{\color{green}\sf\footnotesize\hfill section \ref{#1}}]
               {\color{green}\sf\footnotesize section \ref{#1}\hfill}}
\newcommand{\toExam}[1]
{\marginpar[{\color{green}\sf\footnotesize\hfill example \ref{#1} }]
               {\color{green}\sf\footnotesize example \ref{#1} \hfill}}
\newcommand{\toRem}[1]
{\marginpar[{\color{green}\sf\footnotesize\hfill remark \ref{#1} }]
               {\color{green}\sf\footnotesize remark \ref{#1} \hfill}}
\newcommand{\toPage}[1]
{\marginpar[{\color{green}\sf\footnotesize\hfill p.~\pageref{#1} }]
               {\color{green}\sf\footnotesize p.~\pageref{#1} \hfill}}
%\newcommand{\toCB}[1]
%{\marginpar[{\color{green}\sf\footnotesize\hfill \wwwcb{#1} }]
%               {\color{green}\sf\footnotesize \wwwcb{#1} \hfill}}
\newcommand{\exerbox}[1]
{\marginpar[{\color{red}\sf\footnotesize\hfill exercise \ref{#1} }]
               {\color{red}\sf\footnotesize exercise \ref{#1} \hfill}}
%%%%  mark section as `cyclist':
            \newsavebox{\bike}\sbox{\bike}
            {\raisebox{-5.0ex}[4.5ex][3.5ex]
             {\includegraphics[width=1.6cm]{w11-1bicycle}}} %{Fig/w11-1bicycle}}}
\newcommand{\cyclist}{\marginpar[\hfill\usebox{\bike}]
        {\usebox{\bike}\hfill}}
%%%% mark section as `pedestrian':
           \newsavebox{\trot}\sbox{\trot}
           {\raisebox{-5.0ex}[4.5ex][3.5ex]
            {\includegraphics[width=1.6cm]{i-4-pedestrian}}}
\newcommand{\pedestrian}{\marginpar[\hfill\usebox{\trot}]
        {\usebox{\trot}\hfill}}
\newcommand{\fastTrack}[1]
           {\hfill {\raisebox{-3.0ex}[4.5ex][3.5ex]{
                              \includegraphics[width=1.25cm]{Fig/thisWay.ps}
                                                   }
            \begin{minipage}[t]{0.25\textwidth}
            {\sf \footnotesize
                   fast track: \\
                               \refsect{#1}, p.~\pageref{#1}
                        }
            \end{minipage}
                    }
           }
\newcommand{\fastTrackChap}[1]
           {\hfill {\raisebox{-3.0ex}[4.5ex][3.5ex]{
                              \includegraphics[width=1.25cm]{Fig/thisWay.ps}
                                                   }
            \begin{minipage}[t]{0.25\textwidth}
            {\sf \footnotesize
                   fast track: \\
                               \refchap{#1}, p.~\pageref{#1}
                        }
            \end{minipage}
                    }
           }
\newcommand{\fastTrackExam}[1]
           {{\raisebox{-3.0ex}[-4.5ex][-3.5ex]{
                              \includegraphics[width=1.25cm]{thisWay}
                                                   }
            \ifpaper\else\phantomsection\label{pointer:#1}\fi
            \begin{minipage}[t]{0.16\textwidth}
            {\sf \footnotesize
                   \refexam{#1}\\
                   p.~\pageref{#1}
                        }
            \end{minipage}
                    }
           }
\newcommand{\jumpBack}[1]
{\marginpar[{\color{green}\sf\footnotesize\hfill click to return:~p.~\pageref{pointer:#1} }]
               {\color{green}\sf\footnotesize click to return:~p.~\pageref{pointer:#1} \hfill}}
%\newcommand{\ChapterEnd}{} %\setcounter{endSect}{\thepage}}
        \else
\newcommand{\tangent}{}
\newcommand{\toChap}[1]{}
\newcommand{\toAppe}[1]{}
\newcommand{\toSect}[1]{}
\newcommand{\toExam}[1]{}
\newcommand{\toRem}[1]{}
\newcommand{\toPage}[1]{}
%\newcommand{\toCB}[1]{}
%\newcommand{\exerbox}[1]{}
\newcommand{\cyclist}{}
\newcommand{\pedestrian}{}
\newcommand{\fastTrack}[1]{}
\newcommand{\fastTrackChap}[1]{}
\newcommand{\fastTrackExam}[1]{}
\newcommand{\jumpBack}[1]{}
\newcommand{\ChapterEnd}{}
        \fi

%%%%%%%% for figure layouts:
%% fine tuning floatingfigure \usepackage[nooneline] captions %%%%%
%        \renewcommand{\captionfont}{\footnotesize} % \sffamily}
%        \renewcommand{\captionlabelfont}{\footnotesize \bfseries \rmfamily}
%        % \renewcommand{\captionlabeldelim}{.\quad}
%        % \setlength{\captionmargin}{10pt}

% \FIG{#1}    % \includegraphics[width=0.40\textwidth]{f_name.ps}  ...
%     {#2}    % short caption text omitted here
%     {#3}    % full caption text
%     {#4}    % f-figure-label
%??? 2011-07-23 Predrag: what package defines {widefigure}?
\newcommand{\FIG}[4]{\begin{widefigure}[t]{#4}
    \widefigparts{\begin{minipage}[b]{0.98\textwidth}
                    #1\end{minipage}}{#3}\end{widefigure}}

\newcommand{\FFIG}[5]{\begin{floatingfigure}[v]{0.40\textwidth}
                      \noindent
                      \includegraphics[#1]{#2}
                      \caption[#3]{#4}
                      \label{#5}
              \end{floatingfigure} }

%   \BFIG{#1}   % width=#1\textwidth
%        {#2}   % f_name.ps
%    {#3}   % short caption text
%    {#4}   % full caption text
%    {#5}   % f-figure-label
%       defined here:
\newcommand{\BFIG}[5]{\begin{figure}
              \hspace*{0.10\textwidth}%
              \begin{minipage}[b]{1.00\textwidth}
              \centering{
                      \includegraphics[width=#1\textwidth]{#2}
                     }
                      \caption[#3]{#4}
                      \label{#5}
              \end{minipage}
              \end{figure} }

%  \SFIG{#1}    % f_name.png (or .pdf)
%       {#2}    % short caption text
%       {#3}    % full caption text
%       {#4}    % f-figure-label
\newcommand{\SFIG}[4]{\begin{figure}[h]
              %\hspace*{-0.10\textwidth}
              \hspace*{0.10\textwidth}
              \begin{minipage}[b]{0.55\textwidth}
                      \caption[#2]{#3}
                      \label{#4}
              \end{minipage}~~~~~%
              \begin{minipage}[b]{0.40\textwidth}
                      \includegraphics[width=1.00\textwidth]{#1}
              \end{minipage}
              %\hfill
              \end{figure} }

\newcommand{\MultiFIG}[6]{\begin{figure}[h]
              %\hspace*{-0.10\textwidth}
              \hspace*{0.10\textwidth}
              \begin{minipage}[b]{0.35\textwidth}
                      \caption[#4]{#5}
                      \label{#6}
              \end{minipage}~~~~~%
              \begin{minipage}[b]{0.60\textwidth}
                      \includegraphics[width=1.00\textwidth]{#1}
                      \includegraphics[width=1.00\textwidth]{#2}
                      \includegraphics[width=1.00\textwidth]{#3}
              \end{minipage}
              %\hfill
              \end{figure} }


%%%%%%%%%%%%%%% FORMATING: PAGINATION FOR bin/mkbook %%%%%%%%%%%%%%%%%%%%%%%

\newcommand{\Resume}{\section*{\textsf{\textbf{R\'esum\'e}}}
    %   \addcontentsline{toc}{section}{{~~~~Resum\'e}}
                \addtocontents{toc}{{\small r\'esum\'e \thepage ~}} }
\newcommand{\ResumeEnd}{}

\newcommand{\Remarks}{\section*{\textsf{\textbf{Commentary}}}
%        \addtocontents{toc}{{\small commentary \thepage ~ }}
                     }
\newcommand{\RemarksEnd}{\ifOUP   %\end{enumerate}
        \end{reading} \else  \fi}

\newcommand{\Problems}[2]{
                \renewcommand{\inputfile}{#1 -  #2}
    \ifboyscout \begin{exercisesBS} \else
                \begin{exercises} \fi
    \begin{enumerate}}
\newcommand{\ProblemsEnd}{\end{enumerate}
    \ifboyscout \end{exercisesBS} \else
    \end{exercises} \fi}

\newcommand{\Solution}[4]{
        \section*{Chapter~\ref{c-#1}. #4}
                % PC 2018-01-28, get rid of it in ChaosBook?
                \renewcommand{\inputfile}{#1 -  #2}
                % was \renewcommand{\inputfile}{#1 -  #2}
        % PC:  modified 28 may 2000
                %         \renewcommand{\chapName}{#1}
                %        \setcounter{startSect}{\thepage}
                        }
\newcommand{\Solutions}{
        \section{\textsf{\textbf{Solutions}}}
        %\addcontentsline{toc}{section}{{~~~~Solutions}}
                        }


%%%%%%%%%%%%%%% NUMBERED ENVIRONMENTS %%%%%%%%%%%%%%%%%%%%%%%%%%%%%%%
% REPLACE {13.5cm} by \textwidth counter in these environments

%    \ifarticle
%\newcommand{\exerbox}[1]{}
%\newcommand{\exercise}[2]{}{}
%\newcommand{\solution}[3]{}{}{}
%    \else
%\newtheorem{exerc}{\textsf{\textbf{Exercise}}}[chapter]
%% \newtheorem{exerc}{}[chapter]
%% \newtheorem{exerc}{{$\bullet$}}[chapter]
% \newcommand{\exercise}[2]{
%        %\vskip -13mm
%         \noindent
%         \begin{exerc}{
%\renewcommand{\theenumi}{\alph{enumi}}
%\renewcommand{\labelenumi}{\textbf{(\alph{enumi})\ }}
%    {\noindent\small
%         ~~\textsf{\textbf{#1}} ~
%           \slshape\sffamily{#2}  } % \textsl would not work...
%    }
%         \vskip -1mm
%% removed the line: % \noindent\rule[.1mm]{\linewidth}{.5mm}
%         \end{exerc}
%                          }
%
%\newcommand{\Exercise}[2]{      %environment for obligatory problems
%        %\vskip -13mm
%        \noindent
%        \begin{exerc}{
%\renewcommand{\theenumi}{\alph{enumi}}
%\renewcommand{\labelenumi}
%    {\textsf{\textbf{ (\alph{enumi})\ }}}
%        {\noindent
%         ~~\textsf{\textbf{\underline{#1}}} ~
%           \slshape\sffamily{#2}  } % \textsl would not work...
%        }
%         \vskip -1mm
%        \end{exerc}
%                          }

\renewcommand{\exercise}[2]{
    \item
    {\textbf{#1}} ~ {#2}}

\newcommand{\solution}[3]{
        {\noindent\small
         \textsf{\textbf{Solution \ref{#1}~-~#2}} %LABEL - TITLE
         \slshape\sffamily{#3}                    %TEXT
         }
         \vskip  1ex  %4mm
% removed the line: % \noindent\rule[.1mm]{\linewidth}{.5mm}
                        }
%\newcommand{\exerbox}[1]
%{\marginpar[{\color{red}\sf\footnotesize\hfill exercise \ref{#1} }]
%               {\color{red}\sf\footnotesize exercise \ref{#1} \hfill}}
%    \fi %end of article switch


    \ifarticle
\newtheorem{exmple}{\noindent\small\textsf{\textbf{Example}}}[section]
\newcommand{\example}[2]{
    \vskip -13mm
        %\begin{offset}
        \begin{exmple}
           \noindent\small
           \textsf{\textbf{#1}} ~
       \slshape\sffamily{#2}
       % \textsl would not work...
    \end{exmple}
    %\end{offset}
        \vskip -1mm
             }
    \else
\newtheorem{rmark}{{\small\textsf{\textbf{Remark}}}}[chapter]
\newcommand{\remark}[2]{
        % \begin{quotation}
        \begin{rmark}
        {\small\em\noindent {\small\sf \underline{ #1} ~} #2 }
    \end{rmark}
    % \end{quotation}
              }
\newtheorem{exmple}{\noindent\small\textsf{\textbf{Example}}}[chapter]
\newcommand{\example}[2]{
    \vskip -13mm
        %\begin{offset}
        \begin{exmple}
           \noindent\small
           \textsf{\textbf{#1}} ~
       \slshape\sffamily{#2}
       % \textsl would not work...
    \end{exmple}
    %\end{offset}
        \vskip -1mm
             }
    \fi %end of article switch
